% Generated draft: requires mathematical and source review.
\documentclass[11pt]{article}
\usepackage[margin=25mm]{geometry}
\usepackage{fontspec}
\setmainfont{Latin Modern Roman}
\usepackage{amsmath,amssymb,mathtools,unicode-math}
\setmathfont{Latin Modern Math}
\usepackage{enumitem,needspace,xurl,hyperref}
\setlength{\parindent}{0pt}
\setlength{\parskip}{5pt}
\setlength{\emergencystretch}{3em}
\newcommand{\N}{\mathbb{N}}
\newcommand{\Z}{\mathbb{Z}}
\newcommand{\Q}{\mathbb{Q}}
\newcommand{\R}{\mathbb{R}}
\newcommand{\C}{\mathbb{C}}
\newcommand{\F}{\mathbb{F}}
\newcommand{\A}{\mathbb{A}}
\newcommand{\PP}{\mathbb{P}}
\newcommand{\E}{\mathbb{E}}
\newcommand{\eps}{\varepsilon}
\newcommand{\dd}{\,\mathrm{d}}
\begin{document}
\section*{UnsolvedMath ID: 800410}
\textbf{Source problem code:} catalogue entry
\subsection{Definitions and mathematical statement}
Let \(F=K(B)\) be the function field of an integral algebraic curve \(B\) over a field \(K\). Let \(f:\PP^N_F\to\PP^N_F\) be a morphism of algebraic degree \(d\geq 2\), and let \(P\in\PP^N(F)\). Here algebraic degree means that \(f^*\mathcal O(1)\cong\mathcal O(d)\).

To make the height explicit under the standard geometric function-field interpretation, choose a normal projective curve \(C\) over \(K\) with function field \(F\). Each \(Q\in\PP^N(F)\) extends to a morphism \(q_Q:C\to\PP^N_K\). Define
\[
 h_C(Q)=\deg_K\bigl(q_Q^*\mathcal O(1)\bigr),
 \qquad
 \hat h_f(P)=\lim_{n\to\infty}\frac{h_C(f^{\circ n}(P))}{d^n},
\]
where \(f^{\circ n}\) denotes the \(n\)-fold iterate of \(f\), with \(f^{\circ 0}\) the identity.

\textbf{Question.} Is \(\hat h_f(P)\in\Q\) always true for the function fields, morphisms, and points under consideration? If it is not always true, determine necessary and sufficient conditions on the data \((K(B),f,P)\) for
\[
\hat h_f(P)\in\Q.
\]

\textbf{Unresolved scope qualification.} The supplied excerpt does not specify the characteristic or other hypotheses on \(K\), additional assumptions on \(B\), or the normalization and precise meaning of the function-field height. The definition above records the standard geometric interpretation; its agreement with the intended conventions and scope of the cited question requires review of the source context.
\subsection{Short English statement}
For a morphism of projective space of degree at least two over a curve's function field, ask whether every rational point has rational canonical height. If this fails, characterize exactly when the canonical height is rational. The intended base-field hypotheses and height conventions remain unspecified in the supplied excerpt.
\subsection{Sources}
\begin{itemize}
\item Supplied local UnsolvedMath registry record, permanent ID \(800410\), source code \texttt{catalogue entry}, titled \emph{Open Problems in Arithmetic Dynamics --- Question 33}, in the AMR Open Problem Lists collection.
\item Benedetto, DeMarco, Ingram, Jones et al., \emph{Current Trends and Open Problems in Arithmetic Dynamics} (2018), Question 33. Supplied source link: \url{https://arxiv.org/abs/1806.04980}. The registry attributes the question to these authors and gives the proposed year as 2018.
\item The supplied provenance reports access on 2026-10-16 and extraction by \texttt{arxiv\_latex}. These are supplied provenance details; no independent source visit or verification was performed. The registry labels the problem \texttt{open}; this is a supplied status, not an independently verified literature assessment.
\end{itemize}
\end{document}
